% Part: first-order-logic
% Chapter: introduction
% Section: sentences

\documentclass[../../../include/open-logic-section]{subfiles}

\begin{document}

\olfileid[pt]{fol}{int}{snt}

\olsection{\usetoken{P}{sentence}}

Ok, agora temos uma (esboço de uma) definição de satisfação
(``verdadeira em'') para !!{structure}s e !!{formula}s. Mas isso requer
um ingrediente adicional---uma atribuição de !!{variable}---e o que
queríamos era uma definição de !!{sentence}s. Como nos livramos das
atribuições, e o que são !!{sentence}s?

Você provavelmente se lembra de uma discussão, em sua primeira
introdução à lógica formal, sobre a relação entre !!{variable}s e
quantificadores. Um quantificador é sempre seguido por !!a{variable} e,
então, na parte da !!{sentence} à qual esse quantificador se aplica (seu
``escopo''), entendemos que a !!{variable} está ``ligada'' a esse
quantificador. Em !!{formula}s, não se exigia que toda !!{variable}
tivesse um quantificador correspondente, e as !!{variable}s sem
quantificadores correspondentes são ``livres'' ou ``não ligadas''.
Tomaremos as !!{sentence}s como sendo todas aquelas !!{formula}s que não
têm !!{variable}s livres.

Novamente, a ideia intuitiva de quando uma ocorrência de !!a{variable}
em !!a{formula}~$!A$ está ligada, de qual quantificador a liga e de
quando ela é livre não é difícil de captar. Você pode ter aprendido um
método para testar isso, talvez envolvendo a contagem de parênteses.
Temos que insistir em uma definição precisa---e, como definimos
!!{formula}s por indução, podemos dar uma definição das ocorrências
livres e ligadas de !!a{variable}~$x$ em !!a{formula}~$!A$ também por
indução.  Por exemplo, ela poderia ter este aspecto para nossa
linguagem simplificada:
\begin{enumerate}
  \item Se $!A$ é atômica, todas as ocorrências de $x$ nela são livres
  (isto é, a ocorrência de $x$ em~$\Atom{\Obj P}{x}$ é livre).
  \item Se $!A$ é da forma $\lnot !B$, então uma ocorrência de~$x$
  em~$\lnot !B$ é livre se, e somente se, a ocorrência correspondente
  de~$x$ for livre em~$!B$ (isto é, as ocorrências livres de variáveis
  em~$!B$ são exatamente as ocorrências correspondentes em~$\lnot !B$).
  \item Se $!A$ é da forma $(!B \land !C)$, então uma ocorrência de~$x$
  em~$(!B \land !C)$ é livre se, e somente se, a ocorrência
  correspondente de~$x$ for livre em~$!B$ ou em~$!C$.
  \item Se $!A$ é da forma $\lexists[x][!B]$, então nenhuma ocorrência
  de $x$ em~$!A$ é livre; se é da forma $\lexists[y][!B]$, em que
  $y$ é uma !!{variable} diferente de~$x$, então uma ocorrência de~$x$
  em $\lexists[y][!B]$ é livre se, e somente se, a ocorrência
  correspondente de~$x$ for livre em~$!B$.
\end{enumerate}

Uma vez que tenhamos uma definição precisa de ocorrências livres e
ligadas de variáveis, podemos simplesmente dizer: !!a{sentence} é
qualquer !!{formula} sem ocorrências livres de !!{variable}s.

\end{document}
